\part{Архитектура}
\label{part:architecture}

До этого момента книга в основном смотрела на платёж снаружи:
через участников, протоколы, платёжные системы и операционные
процессы. Но любая команда, которая пытается построить собственный
платёжный продукт или хотя бы всерьёз встроить платежи в свой стек,
рано или поздно упирается в другой вопрос. Как собрать внутреннюю
систему так, чтобы один внешний платёжный поток не распался внутри
на десятки несогласованных состояний, ретраев, таймаутов, валютных
пересчётов, дубликатов и ручных исключений. Иначе говоря, где именно
внутри архитектуры заканчивается \enquote{интеграция со сторонним
провайдером} и начинается настоящая платёжная инженерия.

Пятая часть книги отвечает именно на этот вопрос. Она не пытается
заново объяснить, что такое эквайер, 3-D Secure или chargeback.
Наоборот, она исходит из того, что читатель уже знает внешнюю логику
платежа и теперь хочет понять внутренние границы системы: какой
компонент принимает запрос мерчанта, какой решает, куда отправить
авторизацию, какой хранит токены, какой ведёт обязательства, какой
отвечает за валютную переоценку, какой разговаривает с карточной
сетью, и почему всё это нельзя безнаказанно сложить в один сервис
с одной таблицей транзакций.

Поэтому эта часть построена вокруг внутренних интерфейсов, а не
вокруг внешних брендов. Сначала появляется платёжный шлюз как
пограничный слой между API продукта и платёжной инфраструктурой.
За ним следует внутренний реестр обязательств. Он нужен
не для дублирования банковского учёта, а для фиксации собственного
состояния обязательств и событий внутри продукта.

Дальше разговор естественно расширяется до мультивалютности, где
становится видно, что валюта --- это не просто атрибут суммы, а
отдельный источник временных расхождений, переоценок и проблем
сверки. Затем появляется процессинговый центр как инфраструктурное
ядро, а в финале --- системы эмитента и эквайера, где особенно хорошо
видно, как по-разному устроены две стороны одной и той же транзакции.

В этой части особенно важна мысль о границах ответственности
компонентов. Шлюз не должен притворяться бухгалтерской системой.
Внутренний реестр не должен брать на себя функции сетевого процессинга.
Слой маршрутизации не должен бесконтрольно менять смысл учётных
событий.
Мультивалютный контур не должен жить как косметическая надстройка
над уже проведёнными суммами. Когда эти границы размыты, система может
какое-то время работать, но почти неизбежно начинает терять
идемпотентность, создавать фантомные статусы или запутывать команду
в вопросе о том, где вообще искать источник истины по конкретной
операции.

Читать эту часть полезно как перевод предыдущих разделов книги
в инженерную топологию. Всё, что раньше обсуждалось как внешний
сигнал или операционное последствие, здесь превращается в требования
к компонентам. Fraud и decline management требуют слоёв риска
и маршрутизации. Сверка и диспуты требуют дисциплины событий и
учётных записей. Токенизация и PCI DSS требуют изоляции данных.
Различия между карточными сетями и A2A-контурами требуют архитектуры,
которая умеет жить с несколькими моделями финальности и несколькими
типами внешних статусов одновременно.

Это ещё и часть про зрелость инженерного мышления. На раннем этапе
команда часто спрашивает, какой провайдер подключить. На следующем ---
как повысить долю одобренных операций. Но после определённого масштаба
неизбежно
приходит другой вопрос: какая внутренняя модель позволит пережить
следующие интеграции, новые страны, дополнительные валюты, новые
методы оплаты и растущий объём исключений без постоянного переписывания
ядра. Именно такой вопрос и организует весь дальнейший рассказ.

После него остаётся сделать последний шаг. Даже хорошая архитектура
не спасает, если команда живёт по устаревшим правилам, вторичным
пересказам и случайным PDF из старой папки. Поэтому заключительная
часть \ref{part:navigation} посвящена уже не самому платёжному
механизму, а дисциплине навигации по первоисточникам и изменениям,
без которой ни одна платёжная система не остаётся актуальной надолго.

\begin{kaobox}[title=Три вопроса этой части]
  \begin{itemize}
    \item Как разделить шлюз, внутренний учёт, маршрутизацию
          и процессинг так, чтобы система не путала свои роли.
    \item Почему повторные попытки, валюты, внешние статусы и события требуют
          собственной внутренней модели, а не одной таблицы транзакций.
    \item Где в платёжной архитектуре должны проходить границы
          ответственности компонентов.
  \end{itemize}
\end{kaobox}
